\documentclass[12pt,a4paper]{article}
\usepackage[utf8]{inputenc}
\usepackage[T2A]{fontenc}
\usepackage[russian]{babel}
\usepackage{amsmath,amssymb,amsfonts}
\usepackage{geometry}
\geometry{top=20mm,bottom=20mm,left=15mm,right=15mm}
\usepackage{booktabs}
\usepackage{array}
\usepackage{hyperref}
\usepackage{tcolorbox}

\tcolorboxenvironment{quote}{
    colback=lightgray!15,
    colframe=gray!60,
    boxrule=0.5pt,
    arc=2pt,
    left=8pt,
    right=8pt,
    top=6pt,
    bottom=6pt
}

\title{\textbf{Руководство к контрольной работе № 3 по дисциплине УПОС}\\[0.3em]
\large Кафедра РПУ МЭИ: Воздействие сигнала и шума на тракты РПУ, АД и ЧД}
\author{Справочник формул, физических обоснований и эталонных алгоритмов}
\date{2026 год}

\begin{document}

\maketitle
\tableofcontents
\newpage

\section{Структура приёмного тракта и основные соотношения}

Обобщённая структурная схема супергетеродинного приёмного тракта:
\begin{center}
\textbf{Антенна} ($T_A$) $\xrightarrow{q_1}$ \textbf{ВЦ} ($K_{p.\text{вц}}, t_{\text{вц}}$) $\xrightarrow{q_2}$ \textbf{УРЧ} ($K_{p.\text{урч}}, K_{\text{ш.урч}}$) $\xrightarrow{q_3}$ \textbf{ПЧ} ($K_{p.\text{пч}}, K_{\text{ш.пч}}$) $\xrightarrow{q_4}$ \textbf{УПЧ} ($K_{\text{ш.упч}}, K_{0.\text{упч}}, R_{\text{вх.упч}}$) $\to$ \textbf{Детектор} $\to$ \textbf{БНЧ} ($K_{\text{бнч}}$)
\end{center}

\subsection{Шумовая полоса пропускания $\Pi_{\text{ш}}$}
По условию полоса УПЧ/БВЧ согласована с шириной спектра полезного сигнала:
\begin{itemize}
    \item \textbf{Для АМ-сигнала:} $\Pi_{\text{ш}} = 2 F_{\text{м}}$, где $F_{\text{м}}$ --- максимальная частота модуляции.
    \item \textbf{Для ЧМ-сигнала:} $\Pi_{\text{ш}} = 2 \Delta f_{\text{м}}$, где $\Delta f_{\text{м}}$ --- девиация частоты (при условии широкого спектра $\beta_{\text{чм}} \ge 1$).
\end{itemize}

\subsection{Шум пассивной входной цепи (ВЦ / фидера)}
Входная цепь является пассивным четырёхполюсником с потерями. Коэффициент шума:
\begin{equation}
K_{\text{ш.вц}} = 1 + \frac{T_{\text{вц}}}{T_0} (L_{\text{вц}} - 1) = 1 + \frac{T_{\text{вц}}}{T_0} \left( \frac{1}{K_{p.\text{вц}}} - 1 \right),
\end{equation}
где $T_{\text{вц}} = t_{\text{вц}} [^\circ\text{C}] + 273$ --- абсолютная температура ВЦ в Кельвинах, а $T_0 = 293\text{ К}$ --- стандартная комнатная температура.

\subsection{Формула Фрииса для сквозного коэффициента шума}
При наличии коэффициентов рассогласования стыков ($q_1, q_2, q_3, q_4 \le 1$) проходная мощность уменьшается, что приводит к дополнительному увеличению вклада шумов последующих каскадов:
\begin{equation}
K_{\text{ш.сум}} = K_{\text{ш.вц}} + \frac{K_{\text{ш.урч}} - 1}{K_{p.\text{вц}} \cdot \frac{q_1}{q_2}} + \frac{K_{\text{ш.пч}} - 1}{K_{p.\text{вц}} K_{p.\text{урч}} \cdot \frac{q_1}{q_3}} + \frac{K_{\text{ш.упч}} - 1}{K_{p.\text{вц}} K_{p.\text{урч}} K_{p.\text{пч}} \cdot \frac{q_1}{q_4}}.
\end{equation}
Если коэффициенты рассогласования не заданы или равны $1$, отношение $\frac{q_1}{q_i}$ опускается:
\begin{equation}
K_{\text{ш.сум}} = K_{\text{ш.вц}} + \frac{K_{\text{ш.урч}} - 1}{K_{p.\text{вц}}} + \frac{K_{\text{ш.пч}} - 1}{K_{p.\text{вц}} K_{p.\text{урч}}} + \frac{K_{\text{ш.упч}} - 1}{K_{p.\text{вц}} K_{p.\text{урч}} K_{p.\text{пч}}}.
\end{equation}

\subsection{Шумовые температуры и мощность приведённого шума}
\begin{itemize}
    \item Шумовая температура приёмника (БВЧ):
    \begin{equation}
    T_{\text{ш.бвч}} = T_0 (K_{\text{ш.сум}} - 1).
    \end{equation}
    \item Суммарная эквивалентная шумовая температура на входе РПУ:
    \begin{equation}
    T_{\text{ш}\Sigma} = T_A + T_{\text{ш.бвч}} = T_A + T_0 (K_{\text{ш.сум}} - 1).
    \end{equation}
    \item Мощность шума, приведённого ко входу РПУ:
    \begin{equation}
    P_{\text{ш.вх}} = k \cdot \Pi_{\text{ш}} \cdot T_{\text{ш}\Sigma},
    \end{equation}
    где $k = 1{,}38 \cdot 10^{-23}\text{ Дж/К}$ --- постоянная Больцмана.
\end{itemize}

\newpage
\section{Обязательные словесные формулировки пересчёта сечений}
Преподаватели кафедры требуют текстуального пояснения физических переходов между плоскостями тракта перед использованием расчетных выражений.

\subsection{Переход 1: Эквивалентное приведение шумов ко входу РПУ}
\begin{quote}
\textit{«Для учёта собственных шумов каскадов приёмного тракта воспользуемся методом эквивалентного приведения источников шума ко входу приёмника. Собственная шумовая мощность каждого последующего четырёхполюсника при пересчёте ко входу делится на коэффициенты передачи проходной мощности предшествующих каскадов с учётом коэффициентов рассогласования стыков. Суммарная шумовая температура на входе складывается из шумовой температуры антенны $T_A$ и приведённой шумовой температуры приёмника $T_{\text{ш.бвч}}$».}
\end{quote}

\subsection{Переход 2: Линейность тракта до детектора ($q_{\text{вх.дет}} = q_{\text{вх.рпу}}$)}
\begin{quote}
\textit{«Поскольку линейный тракт приёмника от входа антенны до входа детектора (ВЦ, УРЧ, ПЧ, УПЧ) осуществляет линейное пропорциональное масштабирование уровней сигнала и шума в согласованной полосе пропускания $\Pi_{\text{ш}}$, отношение сигнал-шум по мощности на входе детектора численно равно отношению сигнал-шум на входе приёмника: $q_{\text{вх.дет}} = q_{\text{вх.рпу}}$».}
\end{quote}

\subsection{Переход 3: Физическая мощность шума на входе УПЧ и напряжение на входе АД}
\begin{quote}
\textit{«Приведённая ко входу мощность шума $P_{\text{ш.вх}}$, проходя через преддетекторный тракт до сечения входа УПЧ, масштабируется по сквозному коэффициенту передачи мощности: $P_{\text{ш.вх.упч}} = P_{\text{ш.вх}} \cdot K_{p.\text{вц}} K_{p.\text{урч}} K_{p.\text{пч}}$.\\
Зная активное входное сопротивление УПЧ $R_{\text{вх.упч}}$, находим действующее напряжение шума на его входе: $U_{\text{ш.вх.упч}} = \sqrt{P_{\text{ш.вх.упч}} \cdot R_{\text{вх.упч}}}$.\\
Далее линейный усилитель промежуточной частоты усиливает это напряжение в $K_{0.\text{упч}}$ раз, формируя на входе детектора действующее напряжение шума: $U_{\text{ш.вх.ад}} = K_{0.\text{упч}} \cdot U_{\text{ш.вх.упч}}$».}
\end{quote}

\subsection{Переход 4: Работа амплитудного детектора в сильном поле ($a \gg 1$)}
\begin{quote}
\textit{«Для линейного амплитудного детектора при сильном сигнале ($a \gg 1$) подавление модуляции шумом отсутствует. При 100\%-й гармонической амплитудной модуляции ($m = 1$) отношение сигнал-шум по мощности на выходе детектора равно отношению сигнал-шум на его входе: $q_{\text{вых.ад}} = q_{\text{вх.ад}} = q_{\text{вых}}$.\\
Отношение амплитуды несущей к действующему напряжению шума на входе АД составляет: $a = \frac{U_{c.\text{вх.ад}}}{U_{\text{ш.вх.ад}}} = \sqrt{2 q_{\text{вх.ад}}}$. Так как $a \gg 3$, режим линейного детектирования подтверждён ($M(a) \approx a$)».}
\end{quote}

\subsection{Переход 5: Надпороговый режим работы ЧД}
\begin{quote}
\textit{«В надпороговой области работы частотного детектора спектральная плотность флуктуаций мгновенной частоты шума имеет квадратично-параболический вид: $G_{\text{чд}}(F) \sim F^2$. Интегрирование спектра фильтром БНЧ в полосе $F_{\text{м}}$ обуславливает выигрыш ЧМ: $q_{\text{вых}} = 3 \beta_{\text{чм}}^3 q_{\text{вх.чд}}$.\\
Полученное значение $q_{\text{вх.чд}}$ превосходит пороговое значение $q_{\text{вх.пор}}$ по пороговой кривой для данного индекса $\beta_{\text{чм}}$, следовательно, детектор находится в надпороговом режиме».}
\end{quote}

\newpage
\section{ТИП 1: ЧМ-приёмник. Расчет чувствительности $P_{c.\text{вх min}}$}

\subsection{Условие типовой задачи}
\textbf{Дано:} $\Delta f_{\text{м}} = 600\text{ кГц}$, $F_{\text{м}} = 60\text{ кГц}$, $q_{\text{вых}} = 48\text{ дБ}$, $T_A = 500\text{ К}$, $t_{\text{вц}} = 35^\circ\text{C}$, $K_{p.\text{вц}} = 0{,}82$, $K_{p.\text{урч}} = 3$, $K_{p.\text{пч}} = 0{,}6$, $K_{\text{ш.урч}} = 2$, $K_{\text{ш.пч}} = 3$, $K_{\text{ш.упч}} = 2{,}3$.

\subsection{Пошаговое решение}

\textbf{1. Определение полосы и параметров модуляции:}
\begin{equation}
\Pi_{\text{ш}} = 2 \Delta f_{\text{м}} = 2 \cdot 600 \cdot 10^3 = 1200\text{ кГц} = 1{,}2\text{ МГц},
\end{equation}
\begin{equation}
\beta_{\text{чм}} = \frac{\Delta f_{\text{м}}}{F_{\text{м}}} = \frac{600\text{ кГц}}{60\text{ кГц}} = 10.
\end{equation}

\textbf{2. Расчет требуемого отношения сигнал-шум на входе ЧД:}
Переводим выходное отношение в линейный масштаб:
\begin{equation}
q_{\text{вых}} = 10^{\frac{48}{10}} = 10^{4{,}8} \approx 63095{,}73.
\end{equation}

\textit{Обоснование пересчета:}
\begin{quote}
В надпороговой области работы ЧД энергетический спектр шума на выходе детектора квадратичен ($G_{\text{чд}}(F) \sim F^2$), а срез спектра фильтром БНЧ с граничной частотой $F_{\text{м}}$ даёт выигрыш: $q_{\text{вых}} = 3 \beta_{\text{чм}}^3 q_{\text{вх.чд}}$.
\end{quote}
Отсюда находим:
\begin{equation}
q_{\text{вх.чд}} = \frac{q_{\text{вых}}}{3 \beta_{\text{чм}}^3} = \frac{63095{,}73}{3 \cdot 10^3} = \frac{63095{,}73}{3000} \approx 21{,}03.
\end{equation}
В децибелах:
\begin{equation}
q_{\text{вх.чд [дБ]}} = 10 \lg(21{,}03) \approx 13{,}23\text{ дБ}.
\end{equation}
\textbf{Проверка порогового эффекта:} Для $\beta_{\text{чм}} = 10$ пороговое отношение составляет $q_{\text{вх.пор}} \approx 10\dots11\text{ дБ}$. Так как $13{,}23\text{ дБ} > q_{\text{вх.пор}}$, детектор гарантированно работает в надпороговой области, и формула $3 \beta_{\text{чм}}^3$ корректна.

\textbf{3. Определение сквозного шума приёмного тракта:}
\textit{Обоснование приведения ко входу:}
\begin{quote}
Для учёта собственных шумов блоков воспользуемся методом приведения шумов ко входу РПУ через формулу Фрииса.
\end{quote}
Температура ВЦ в Кельвинах: $T_{\text{вц}} = 273 + 35 = 308\text{ К}$.
\begin{equation}
K_{\text{ш.вц}} = 1 + \frac{308}{293} \left( \frac{1}{0{,}82} - 1 \right) \approx 1{,}231.
\end{equation}
Сквозной коэффициент шума тракта:
\begin{equation}
K_{\text{ш.сум}} = 1{,}231 + \frac{2 - 1}{0{,}82} + \frac{3 - 1}{0{,}82 \cdot 3} + \frac{2{,}3 - 1}{0{,}82 \cdot 3 \cdot 0{,}6} \approx 1{,}231 + 1{,}220 + 0{,}813 + 0{,}881 = 4{,}145.
\end{equation}
Шумовая температура приёмника:
\begin{equation}
T_{\text{ш.бвч}} = T_0 (K_{\text{ш.сум}} - 1) = 293 \cdot (4{,}145 - 1) \approx 921{,}48\text{ К}.
\end{equation}
Суммарная шумовая температура тракта:
\begin{equation}
T_{\text{ш}\Sigma} = T_A + T_{\text{ш.бвч}} = 500 + 921{,}48 = 1421{,}48\text{ К}.
\end{equation}
Мощность шума, приведённого ко входу РПУ:
\begin{equation}
P_{\text{ш.вх}} = k \cdot \Pi_{\text{ш}} \cdot T_{\text{ш}\Sigma} = 1{,}38 \cdot 10^{-23} \cdot 1{,}2 \cdot 10^6 \cdot 1421{,}48 \approx 2{,}354 \cdot 10^{-14}\text{ Вт}.
\end{equation}

\textbf{4. Определение чувствительности:}
\textit{Обоснование пересчета отношения с/ш:}
\begin{quote}
Линейный преддетекторный тракт сохраняет отношение сигнал-шум по мощности неизменным: $q_{\text{вх.рпу}} = q_{\text{вх.чд}} \approx 21{,}03$.
\end{quote}
\begin{equation}
P_{c.\text{вх min}} = q_{\text{вх.рпу}} \cdot P_{\text{ш.вх}} = 21{,}03 \cdot 2{,}354 \cdot 10^{-14} \approx 4{,}95 \cdot 10^{-13}\text{ Вт}.
\end{equation}

\textbf{Ответ:} $P_{c.\text{вх min}} \approx 4{,}95 \cdot 10^{-13}\text{ Вт}$.

\newpage
\section{ТИП 2: ЧМ-приёмник. Обратная задача (нахождение $t_{\text{вц max}}$)}

\subsection{Условие типовой задачи}
\textbf{Дано:} На входе действует сигнал мощностью $P_c = 1{,}24 \cdot 10^{-13}\text{ Вт}$, $\Delta f_{\text{м}} = 125\text{ кГц}$, $F_{\text{м}} = 25\text{ кГц}$, требуемое $q_{\text{вых}} = 40\text{ дБ}$, $T_A = 200\text{ К}$, $K_{p.\text{вц}} = 0{,}72$, $K_{p.\text{урч}} = 9$, $K_{p.\text{пч}} = 0{,}6$, $K_{\text{ш.урч}} = 1{,}5$, $K_{\text{ш.пч}} = 2$, $K_{\text{ш.упч}} = 2{,}5$. Все $q_i = 1$.\\
\textbf{Найти:} Максимально допустимую физическую температуру входной цепи $t_{\text{вц max}}$ ($^\circ\text{C}$).

\subsection{Пошаговое решение}

\textbf{1. Параметры сигнала:}
\begin{equation}
\Pi_{\text{ш}} = 2 \Delta f_{\text{м}} = 250\text{ кГц}, \quad \beta_{\text{чм}} = \frac{125}{25} = 5.
\end{equation}
Выходное отношение сигнал-шум в разах:
\begin{equation}
q_{\text{вых}} = 10^{\frac{40}{10}} = 10^4.
\end{equation}

\textbf{2. Отношение сигнал-шум на входе ЧД:}
\textit{Обоснование формулы ЧД:}
\begin{quote}
Предполагая надпороговую работу ЧД, выходное отношение с/ш формируется в результате выигрыша ЧМ: $q_{\text{вых}} = 3 \beta_{\text{чм}}^3 q_{\text{вх.чд}}$.
\end{quote}
\begin{equation}
q_{\text{вх.чд}} = \frac{q_{\text{вых}}}{3 \beta_{\text{чм}}^3} = \frac{10^4}{3 \cdot 5^3} = \frac{10^4}{375} \approx 26{,}67 \quad (\approx 14{,}26\text{ дБ}).
\end{equation}
Для $\beta_{\text{чм}} = 5$ порог $q_{\text{вх.пор}} \approx 9\dots10\text{ дБ}$. Режим строго надпороговый ($14{,}26 > 10$).

\textbf{3. Предельно допустимая мощность шума на входе:}
\textit{Обоснование пересчета сечения:}
\begin{quote}
Так как линейный тракт до детектора масштабирует сигнал и шум в равной степени, отношение сигнал-шум на входе приёмника совпадает со входным отношением детектора: $q_{\text{вх.рпу}} = q_{\text{вх.чд}}$. Отсюда мощность приведённого входного шума не должна превышать пороговую:
\end{quote}
\begin{equation}
P_{\text{ш.вх max}} = \frac{P_c}{q_{\text{вх.рпу}}} = \frac{1{,}24 \cdot 10^{-13}\text{ Вт}}{26{,}67} \approx 4{,}65 \cdot 10^{-15}\text{ Вт}.
\end{equation}

\textbf{4. Максимально допустимая шумовая температура тракта:}
\begin{equation}
T_{\text{ш}\Sigma \max} = \frac{P_{\text{ш.вх max}}}{k \cdot \Pi_{\text{ш}}} = \frac{4{,}65 \cdot 10^{-15}}{1{,}38 \cdot 10^{-23} \cdot 250 \cdot 10^3} \approx 1347{,}8\text{ К}.
\end{equation}
\textit{Обоснование вычитания шума антенны:}
\begin{quote}
Суммарная шумовая температура формируется аддитивным сложением шума антенны и собственных приведённых шумов приемного тракта: $T_{\text{ш}\Sigma} = T_A + T_{\text{ш.бвч}}$.
\end{quote}
\begin{equation}
T_{\text{ш.бвч max}} = T_{\text{ш}\Sigma \max} - T_A = 1347{,}8 - 200 = 1147{,}8\text{ К}.
\end{equation}

\textbf{5. Предельный суммарный коэффициент шума:}
\begin{equation}
K_{\text{ш.сум max}} = 1 + \frac{T_{\text{ш.бвч max}}}{T_0} = 1 + \frac{1147{,}8}{293} \approx 4{,}918.
\end{equation}

\textbf{6. Расчёт шума ВЦ по формуле Фрииса:}
\textit{Обоснование вычитания шумов активных каскадов:}
\begin{quote}
Шумы активных каскадов (УРЧ, ПЧ, УПЧ) фиксированы параметрами схемы. Вычитая их приведённый вклад из суммарного допустимого коэффициента шума, находим максимально допустимый коэффициент шума входной цепи $K_{\text{ш.вц max}}$:
\end{quote}
\begin{equation}
K_{\text{ш.вц max}} = K_{\text{ш.сум max}} - \left( \frac{K_{\text{ш.урч}} - 1}{K_{p.\text{вц}}} + \frac{K_{\text{ш.пч}} - 1}{K_{p.\text{вц}} K_{p.\text{урч}}} + \frac{K_{\text{ш.упч}} - 1}{K_{p.\text{вц}} K_{p.\text{урч}} K_{p.\text{пч}}} \right).
\end{equation}
Подставляем численные значения:
\begin{equation}
\Delta K_{\text{ш.тракт}} = \frac{1{,}5 - 1}{0{,}72} + \frac{2 - 1}{0{,}72 \cdot 9} + \frac{2{,}5 - 1}{0{,}72 \cdot 9 \cdot 0{,}6} \approx 0{,}6944 + 0{,}1543 + 0{,}3858 \approx 1{,}2345.
\end{equation}
\begin{equation}
K_{\text{ш.вц max}} = 4{,}918 - 1{,}2345 \approx 3{,}6835.
\end{equation}

\textbf{7. Нахождение температуры ВЦ:}
Используем зависимость коэффициента шума пассивного четырёхполюсника от температуры:
\begin{equation}
K_{\text{ш.вц}} = 1 + \frac{T_{\text{вц}}}{T_0} \left( \frac{1}{K_{p.\text{вц}}} - 1 \right) \implies T_{\text{вц max}} = \frac{(K_{\text{ш.вц max}} - 1) T_0}{\frac{1}{K_{p.\text{вц}}} - 1}.
\end{equation}
\begin{equation}
T_{\text{вц max}} = \frac{(3{,}6835 - 1) \cdot 293}{\frac{1}{0{,}72} - 1} = \frac{2{,}6835 \cdot 293}{0{,}3889} \approx 2021{,}7\text{ К}.
\end{equation}
Перевод в градусы Цельсия:
\begin{equation}
t_{\text{вц max}} = T_{\text{вц max}} - 273 \approx 2021{,}7 - 273 \approx 1748{,}7^\circ\text{C}.
\end{equation}

\textbf{Ответ:} $t_{\text{вц max}} \approx 1749^\circ\text{C}$.

\newpage
\section{ТИП 3: АМ-приёмник. Расчет амплитуды $U_{c.\text{вых}}$}

\subsection{Условие типовой задачи}
\textbf{Дано:} Гармоническая 100\%-я АМ ($m = 1$), $F_{\text{м}} = 100\text{ кГц}$, $q_{\text{вых}} = 20\text{ дБ}$, $T_A = 500\text{ К}$, $t_{\text{вц}} = 40^\circ\text{C}$, $K_{p.\text{вц}} = 0{,}8$, $K_{p.\text{урч}} = 3{,}5$, $K_{p.\text{пч}} = 0{,}6$, $K_{0.\text{упч}} = 1000$, $R_{\text{вх.упч}} = 1\text{ кОм}$, $K_{\text{д}} = 0{,}8$, $K_{\text{бнч}} = 20$, $K_{\text{ш.урч}} = 2$, $K_{\text{ш.пч}} = 3$, $K_{\text{ш.упч}} = 2{,}3$.

\subsection{Пошаговое решение}

\textbf{1. Полоса и шум на входе РПУ:}
Полоса УПЧ равна ширине спектра АМ-колебания:
\begin{equation}
\Pi_{\text{ш}} = 2 F_{\text{м}} = 2 \cdot 100 \cdot 10^3 = 200\text{ кГц}.
\end{equation}
Температура ВЦ: $T_{\text{вц}} = 273 + 40 = 313\text{ К}$.
\begin{equation}
K_{\text{ш.вц}} = 1 + \frac{313}{293} \left( \frac{1}{0{,}8} - 1 \right) \approx 1{,}267.
\end{equation}
Суммарный шум БВЧ:
\begin{equation}
K_{\text{ш.сум}} = 1{,}267 + \frac{2 - 1}{0{,}8} + \frac{3 - 1}{0{,}8 \cdot 3{,}5} + \frac{2{,}3 - 1}{0{,}8 \cdot 3{,}5 \cdot 0{,}6} \approx 1{,}267 + 1{,}250 + 0{,}714 + 0{,}774 = 4{,}005.
\end{equation}
\begin{equation}
T_{\text{ш.бвч}} = 293 \cdot (4{,}005 - 1) \approx 880{,}5\text{ К}, \quad T_{\text{ш}\Sigma} = 500 + 880{,}5 = 1380{,}5\text{ К}.
\end{equation}
Мощность шума на входе антенны:
\begin{equation}
P_{\text{ш.вх}} = 1{,}38 \cdot 10^{-23} \cdot 200 \cdot 10^3 \cdot 1380{,}5 \approx 3{,}81 \cdot 10^{-15}\text{ Вт}.
\end{equation}

\textbf{2. Пересчёт шума к сечению входа АД:}
\textit{Обоснование масштабирования мощности шума до входа УПЧ:}
\begin{quote}
Мощность приведённого входного шума, проходя через преддетекторные каскады (ВЦ, УРЧ, ПЧ), масштабируется на произведение их сквозных коэффициентов передачи проходной мощности:
\end{quote}
\begin{equation}
P_{\text{ш.вх.упч}} = P_{\text{ш.вх}} \cdot K_{p.\text{вц}} K_{p.\text{урч}} K_{p.\text{пч}} = 3{,}81 \cdot 10^{-15} \cdot 0{,}8 \cdot 3{,}5 \cdot 0{,}6 \approx 6{,}40 \cdot 10^{-15}\text{ Вт}.
\end{equation}
\textit{Обоснование пересчета мощности в напряжение:}
\begin{quote}
Зная входное активное сопротивление УПЧ $R_{\text{вх.упч}} = 1000\text{ Ом}$, находим эффективное напряжение шума на его входе, а затем умножаем на коэффициент усиления УПЧ по напряжению $K_{0.\text{упч}}$:
\end{quote}
\begin{equation}
U_{\text{ш.вх.упч}} = \sqrt{P_{\text{ш.вх.упч}} \cdot R_{\text{вх.упч}}} = \sqrt{6{,}40 \cdot 10^{-15} \cdot 1000} \approx 2{,}53 \cdot 10^{-6}\text{ В}.
\end{equation}
\begin{equation}
U_{\text{ш.вх.ад}} = K_{0.\text{упч}} \cdot U_{\text{ш.вх.упч}} = 1000 \cdot 2{,}53 \cdot 10^{-6} = 2{,}53 \cdot 10^{-3}\text{ В}.
\end{equation}

\textbf{3. Определение амплитуды несущей на входе АД:}
Выходное отношение сигнал-шум в разах:
\begin{equation}
q_{\text{вых}} = 10^{\frac{20}{10}} = 100.
\end{equation}
\textit{Обоснование режима работы детектора:}
\begin{quote}
Для линейного АД в сильном поле при $m=1$ отношение с/ш на входе детектора совпадает с выходным: $q_{\text{вх.ад}} = q_{\text{вых}} = 100$. Отношение амплитуды несущей к эффективному напряжению шума определяется формулой $a = \sqrt{2 q_{\text{вх.ад}}}$.
\end{quote}
\begin{equation}
a = \sqrt{2 \cdot 100} \approx 14{,}14.
\end{equation}
Условие $a \gg 3$ выполняется, режим строго линейный ($M(a) \approx a$).
Амплитуда несущей на входе АД:
\begin{equation}
U_{c.\text{вх.ад}} = a \cdot U_{\text{ш.вх.ад}} = 14{,}14 \cdot 2{,}53 \cdot 10^{-3} \approx 3{,}58 \cdot 10^{-2}\text{ В}.
\end{equation}

\textbf{4. Определение выходного демодулированного напряжения:}
\textit{Обоснование детектирования и фильтрации БНЧ:}
\begin{quote}
Линейный диодный детектор выделяет огибающую с коэффициентом передачи $K_{\text{д}}$. При 100\% модуляции амплитуда низкочастотной составляющей на выходе детектора равна $U_{\text{д}} = K_{\text{д}} \cdot m \cdot U_{c.\text{вх.ад}}$. БНЧ усиливает этот сигнал в $K_{\text{бнч}}$ раз.
\end{quote}
\begin{equation}
U_{\text{д}} = 0{,}8 \cdot 1 \cdot 3{,}58 \cdot 10^{-2} \approx 2{,}864 \cdot 10^{-2}\text{ В},
\end{equation}
\begin{equation}
U_{c.\text{вых}} = K_{\text{бнч}} \cdot U_{\text{д}} = 20 \cdot 2{,}864 \cdot 10^{-2} \approx 0{,}573\text{ В}.
\end{equation}

\textbf{Ответ:} $U_{c.\text{вых}} \approx 0{,}573\text{ В}$.

\newpage
\section{ТИП 4: Спектральная плотность и воздействие параметров цепей на ЧД}

\subsection{Энергетический спектр на выходе ЧД и напряжение шума на выходе БНЧ}
В надпороговой области ($a \gg 1$) энергетический спектр шума на выходе ЧД описывается выражением:
\begin{equation}
G_{\text{чд}}(F) = S_{\text{чд}}^2 \frac{2 G_{\text{ш}}}{U_c^2} F^2, \quad F \in \left[0, \frac{\Pi_{\text{ш}}}{2}\right],
\end{equation}
где $S_{\text{чд}}$ --- крутизна ЧД (переводится в $\text{В/Гц}$), $G_{\text{ш}} = \frac{U_{\text{ш}}^2}{\Pi_{\text{ш}}}$ --- спектральная плотность дисперсии шума на входе ЧД, $U_c$ --- амплитуда сигнала на входе ЧД.

Дисперсия и эффективное напряжение шума на выходе БНЧ с прямоугольной полосой $F_{\text{м}} \le \frac{\Pi_{\text{ш}}}{2}$:
\begin{equation}
U_{\text{ш.вых}}^2 = \int_0^{F_{\text{м}}} K_{\text{бнч}}^2 G_{\text{чд}}(F) dF = K_{\text{бнч}}^2 S_{\text{чд}}^2 \frac{2 G_{\text{ш}}}{U_c^2} \int_0^{F_{\text{м}}} F^2 dF = K_{\text{бнч}}^2 S_{\text{чд}}^2 \frac{2 G_{\text{ш}} F_{\text{м}}^3}{3 U_c^2}.
\end{equation}
\begin{equation}
U_{\text{ш.вых}} = K_{\text{бнч}} S_{\text{чд}} \sqrt{\frac{2 G_{\text{ш}} F_{\text{м}}^3}{3 U_c^2}}.
\end{equation}

\subsection{Теоретический вопрос: Расширение полосы БВЧ}
\textbf{Задача:} Во сколько раз изменится мощность шума на выходе БНЧ, если полосу пропускания БВЧ $\Pi_{\text{ш}}$ расширить в $n$ раз (например, в $3$ раза), сохранив коэффициент усиления тракта и амплитуду немодулированного сигнала $U_c$?\\
\textbf{Решение и словесное обоснование:}
\begin{quote}
Формула для мощности шума на выходе БНЧ зависит только от спектральной плотности белого шума на входе ЧД $G_{\text{ш}}$, крутизны $S_{\text{чд}}$, амплитуды сигнала $U_c$ и полосы фильтра БНЧ $F_{\text{м}}$. Сама полоса пропускания БВЧ $\Pi_{\text{ш}}$ определяет только верхнюю границу существования квадратичного спектра $\frac{\Pi_{\text{ш}}}{2}$.\\
Поскольку после расширения полосы граничная частота БНЧ по-прежнему удовлетворяет условию $F_{\text{м}} \le \frac{n \Pi_{\text{ш}}}{2}$, пределы интегрирования спектра не меняются.
\end{quote}
\textbf{Проверка надпорогового условия:}
При увеличении полосы БВЧ в $n$ раз общая дисперсия шума на входе детектора увеличивается в $n$ раз: $U_{\text{ш}}' = U_{\text{ш}} \sqrt{n}$.
Соответственно, отношение амплитуд уменьшается:
\begin{equation}
a' = \frac{U_c}{U_{\text{ш}}'} = \frac{a}{\sqrt{n}}, \quad q_{\text{вх.чд}}' = \frac{(a')^2}{2} = \frac{q_{\text{вх.чд}}}{n}.
\end{equation}
Если полученное значение $q_{\text{вх.чд}}' > q_{\text{вх.пор}}$, детектор остаётся в надпороговой области.\\
\textbf{Вывод:} Мощность шума на выходе БНЧ \textbf{не изменится} (отношение мощностей равно $1$).

\newpage
\section{ТИП 5: ЧД в подпороговой области (Аномальный режим)}
Если при расчете чувствительности или спектра оказывается, что $q_{\text{вх.чд}} < q_{\text{вх.пор}}$ (например, $q_{\text{вх.чд}} \approx 4\dots 6\text{ дБ}$ при пороге $10\text{ дБ}$), формула $q_{\text{вых}} = 3 \beta_{\text{чм}}^3 q_{\text{вх.чд}}$ даёт завышенный ошибочный результат.

\subsection{Физическая природа подпорогового эффекта}
\begin{quote}
\textit{«При малом отношении сигнал-шум ($a \approx 1\dots 3$) годограф комплексной огибающей результирующего колебания охватывает начало координат, что вызывает скачки полной фазы на $2\pi$. Дифференцирование этих скачков в частотном детекторе порождает мощные кратковременные импульсы мгновенной частоты (аномальные выбросы Райса). Их спектр равномерен в области низких частот, что резко увеличивает шум на выходе БНЧ и вызывает скачкообразное падение выходного отношения сигнал-шум».}
\end{quote}

\subsection{Формула Райса для подпорогового режима}
Отношение сигнал-шум на выходе БНЧ с учётом импульсного шума выражается формулой:
\begin{equation}
q_{\text{вых}} = \frac{3 \beta_{\text{чм}}^3 q_{\text{вх.чд}}}{1 + \frac{24}{\pi} \beta_{\text{чм}}^2 q_{\text{вх.чд}} e^{-q_{\text{вх.чд}}}}.
\end{equation}
\begin{itemize}
    \item В надпороговой области $q_{\text{вх.чд}} \gg 1 \implies e^{-q_{\text{вх.чд}}} \to 0$, и знаменатель равен $1$ (переход в стандартную формулу $3 \beta_{\text{чм}}^3 q_{\text{вх}}$).
    \item В подпороговой области знаменатель больше единицы, что описывает резкую деградацию помехоустойчивости. Точное значение $q_{\text{вых}}$ снимается по нелинейным кривым моделей, представленным в раздаточных материалах кафедры.
\end{itemize}

\end{document}